
\section{Conclusions}
In this paper we presented an extension of the approach of \cite{article_KRPT2009} to the more general setting of affine processes on positive definite matrices. We showed that their methodology may be adapted to this state space in a straightforward way. What is more, it is possible to efficiently price European swaptions in this multi-factor setting by means of a cumulant expansion due to \cite{article_CollinDufresneGoldstein2002}. In doing so we are in front of a setting which is potentially able to capture correlation effects which can not be described by a single-factor framework. We provided numerical examples for the Wishart Libor model, where the introduction of off-diagonal elements gives rise to new possibilities in the control of the shape of the implied volatility surface.\\

Our contribution may be seen as a starting point for a description of market models in this state space, in consequence we believe that there are many possible directions for future research. An example is given by the problem of calibrating this family of models to real market data. As the structure of the products in the fixed-income market suggests, even in the plain vanilla case, we expect the objective function that should be minimized in the calibration procedure to be quite involved. Yet, some calibration results were obtained on equity derivatives in \cite{dafgra11} for Wishart based models so a calibration using interest rates derivatives might be feasible. Certainly, it will be a delicate issue and may constitute an interesting contribution by its own. Once the model is calibrated on vanillas, one could then further investigate the performance of the model on more exotic structures, like e.g. Bermudan swaptions and barrier options. Theses issues are left for future work.

